\chapter{Decode a PNG through several representations}
\MARSChapterHeading{Decode a PNG through several representations}
\label{revision:png}

A PNG file does not contain the display words that the bitmap tool reads. Its
chunks contain a compressed stream; that stream expands into filtered rows;
those rows must be reconstructed before their red, green and blue samples can
be packed into words. A plausible result at one boundary does not establish
that the next representation is valid.

The complete \texttt{png-viewer.asm} in \texttt{revision/labs/files} joins these
stages. This chapter follows a two-by-five image through the program, explains
the state required for DEFLATE, and derives the dependencies that permit some
vector operations while preventing others. It assumes the earlier CRC, bounded
file, QOI and RGB/MSA work. The lab retains the full decoder construction route
and assigns a forward row filter independently.

Our local profile accepts noninterlaced RGB PNGs with eight bits per channel,
width and height from one through 512, and at most 1 MiB of file bytes. It
supports stored, fixed-Huffman and dynamic-Huffman DEFLATE blocks and all five
PNG row filters. It rejects other color formats, transparency and animation.
It preserves sample numbers without color management. Those restrictions define
this exercise, rather than all valid PNG files.

\section{Give each stage a different count}

The trace fixture contains ten pixels. Its complete file is 115 bytes, its
joined compressed stream is 46 bytes, and its inflated rows occupy 35 bytes.
The RGB samples account for thirty of those bytes; five more select the row
filters. The packed display pixels occupy forty active bytes.

\begin{table}[!ht]
  \centering
  \begin{tabular}{l r l}
    \toprule
    Representation        & Trace bytes & Size for width $w$, height $h$  \\
    \midrule
    Complete PNG          & 115         & Bounded supplied file count     \\
    Joined IDAT payloads  & 46          & Sum of consecutive IDAT lengths \\
    Inflated row buffer   & 35          & $h(1+3w)$                       \\
    RGB sample bytes      & 30          & $3wh$                           \\
    Active display pixels & 40          & $4wh$                           \\
    \bottomrule
  \end{tabular}
  \caption{Each boundary establishes a different representation and count.}
\end{table}

The ordinary framebuffer has a stride of 512 words: successive image rows
start 2048 bytes apart. Its forty active bytes are therefore spread across
five rows. With a tightly packed stride of two words, they would be adjacent.
Stride changes destination placement; it does not change the compressed stream
or the image's pixel count.

The caller reserves separate file, compressed, row and framebuffer storage.
Its context holds stride and framebuffer capacity, compressed address and
capacity, then row address and capacity, as six native doublewords at offsets
zero through forty. The final three doublewords at offsets 48, 56 and 64 are
width, height and active pixel-byte results. Native fields follow the selected
MARS byte order. File fields retain the byte order defined by their format.

\section{Collect one stream from checked chunks}

After the eight-byte PNG signature, each chunk has a big-endian unsigned
four-byte payload length, four type bytes, its payload and a four-byte CRC.
There is no odd-byte padding. For header offset $q$ and payload length $L$,
the next header is at $q+12+L$. Prove the envelope fits before reading its
payload or stored CRC. Decode the length unsigned and perform offset arithmetic
at address width.

In the trace, IHDR begins at offset eight, has thirteen payload bytes, and
ends at 33. Two IDAT chunks start at 33 and 46, holding one and 45 payload
bytes respectively. IEND starts at 103 and ends at 115. The first IDAT contains
only \texttt{78}; the next begins \texttt{01}. Even the two-byte zlib header
crosses this chunk boundary. A fresh inflater for each IDAT would be wrong.
Concatenate their payloads, excluding chunk envelopes, to obtain one stream.

Require IHDR first, one uninterrupted IDAT run, and an empty IEND at the exact
file end. Check every chunk's CRC over its type and payload, including an
unknown ancillary chunk that will be skipped. Length bytes are excluded from
that CRC. Validate four ASCII letters in each type and an uppercase reserved
third letter. Unknown critical chunks fail. Our profile permits an optional
valid PLTE before IDAT, but truecolor samples do not index it. It explicitly
rejects \texttt{tRNS} and animation chunks; ignoring those could display a
misleading image. An ancillary chunk after IDAT closes that run.

\texttt{png\_collect(file,count,compressed,descriptor)} takes destination
capacity in the descriptor's first native doubleword. On success, it writes
width, height, joined length and predicted inflated length into the following
four doublewords, and returns zero with joined length in \texttt{v1}.
On failure it returns one and zero and retains the whole descriptor. Compressed
scratch may already contain a copied prefix. The prediction $h(1+3w)$ derives
from validated dimensions; it is not evidence that the stream will produce
those bytes. Collection deliberately leaves that question to the inflater.

\section{Separate the wrapper, bits and symbols}

The joined stream begins with a zlib header and ends with a big-endian Adler-32
checksum. Validate DEFLATE as the method, the advertised window, the header's
check bits and absence of a preset dictionary. Reserve the four trailer bytes
from the compressed-bit input window. A bit refill must not borrow checksum
bytes to finish a truncated block.

The trace stream begins \texttt{78 01 01 23 00 dc ff}. The first two bytes
are the wrapper. The next byte begins a final stored block: its low bit marks
the final block, its next two bits select stored data, and the remaining bits
are discarded to reach a byte boundary. The little-endian stored length is
\texttt{23 00}, or 35. Its complement is \texttt{dc ff}. Exactly 35 row bytes
follow, then Adler-32. Even within this one stream, the wrapper's multibyte
fields and stored-block length have different byte orders.

A bit reservoir keeps unread bits and a separate available-bit count. To read
an ordinary $n$-bit field, refill with bounded byte loads, take the low $n$
bits, shift right by $n$, and decrease the count. Block type one uses a fixed
Huffman alphabet; type two transmits code lengths from which alphabets are
built; type three is invalid. Nonfinal blocks continue into another block
header without resetting produced history.

Huffman codes require a separate interpretation. Suppose an alphabet has codes
\texttt{0}, \texttt{10} and \texttt{11}. Arriving bits one then zero select
code \texttt{10}. Interpreting those arrivals as an ordinary low-bit-first
two-bit number gives one. A symbol decoder can consume one bit at a time,
append it to the developing code, and compare against the canonical codes of
that length. Keep code numbering separate from the reservoir's numerical
field extraction.

A dynamic block first describes the alphabet used to encode its own code
lengths, then the literal/length and distance alphabets. Length zero means a
symbol is absent. Repeat operations expand earlier lengths or runs of zeros;
prove each run fits the remaining length array before storing it. A repeat of
the previous length is invalid before any previous length exists. Clear and
rebuild each block's tables so a preceding block cannot supply missing entries.

\Needspace{13\baselineskip}
\begin{table}[!ht]
  \centering
  \begin{tabular}{r l r l}
    \toprule
    Symbol & Meaning      & Code length & Canonical code \\
    \midrule
    65     & Literal A    & 1           & \texttt{0}     \\
    256    & End of block & 2           & \texttt{10}    \\
    257    & Length three & 2           & \texttt{11}    \\
    \bottomrule
  \end{tabular}
  \caption{A small complete literal/length alphabet. Other symbols have length zero.}
\end{table}

Canonical assignment orders shorter codes first and equal-length codes by
symbol number. After the one-bit code zero, the next available prefix shifts
left to binary ten; symbols 256 and 257 receive ten and eleven. Count available
prefixes while building the table: too many codes at any length is an error.
Incomplete alphabets also need explicit validation, including the permitted
single-symbol and empty-distance cases. The lab develops those rules before
asking you to process arbitrary dynamic headers.

\section{Expand matches into bytes that really exist}

A literal emits one byte. A length/distance pair copies earlier output. In the
small alphabet above, a literal A followed by length three and distance one
produces four A bytes. End-of-block then finishes that compressed block.
A distance of zero, a distance beyond produced history or the advertised window,
a reserved symbol, or a match exceeding remaining output space must fail before
the corresponding load or store.

A match can overlap itself. Starting with one A, length six and distance one
appends six more A bytes. The byte just written supplies the next iteration.
The worked copy loop begins only after distance and output bounds have been
established. Register \texttt{\$t0} is the history source, \texttt{\$s4} the
next output address and \texttt{\$t1} the required output end:
\Needspace{11\baselineskip}
\MARSCodeLabel{Allow newly produced bytes to extend a match}
\begin{minted}{mars-idea}
    lbu $t2, 0($t0)              # forward copy also handles overlapping matches
    sb $t2, 0($s4)
    daddiu $t0, $t0, 1
    daddiu $s4, $s4, 1
    bnec $s4, $t1, copy_match
    nop
\end{minted}
Preloading the whole match in a vector could read bytes that have not yet been
produced. A backward overlap copy would also change this expansion rule.
Optimize only after distinguishing matches whose inputs already exist from
those that extend their own history.

\texttt{png\_inflate(stream,count,rows,expected)} requires the exact expected
output count, rather than merely a capacity ceiling. On success it returns zero
and that produced count. A short stream cannot succeed by leaving old workspace
bytes in place, even if a checksum collision makes the final Adler value agree.
Adler-32 starts with $s_1=1,s_2=0$ and updates $s_1$ with each produced byte,
then $s_2$ with the new $s_1$, modulo 65521. Count only bytes actually emitted.
Require final block completion, the exact output end, a complete wrapper trailer
and no extra compressed bytes under this interface's exact-stream contract.
Failure returns one and zero; selected row scratch may have changed.

\section{Reconstruct colors using reconstructed neighbors}

Inflated data is still not RGB. Every row starts with a filter selector and
then $3w$ residual bytes. For each sample byte, call the reconstructed sample
of the same channel in the previous pixel $a$, the one above $b$, and the upper
left one $c$. Missing neighbors are zero. The channel spacing is three bytes,
so the left red sample is three bytes earlier, not the immediately preceding
blue sample.

The five predictors are zero for None, $a$ for Sub, $b$ for Up,
$\lfloor(a+b)/2\rfloor$ for Average, and a selected neighbor for Paeth.
Paeth forms $p=a+b-c$ and chooses the neighbor closest to $p$, breaking ties
in the order $a$, then $b$, then $c$. Compute differences and the Average sum
without byte wrapping. Reconstruction stores
$(residual+prediction)\bmod256$. This wrapping restores an unsigned sample;
it is not the audio saturation rule.

\Needspace{16\baselineskip}
\begin{table}[!ht]
  \centering
  \begin{tabular}{r l l l}
    \toprule
    Row & Filter  & Residuals, decimal & Reconstructed RGB bytes \\
    \midrule
    0   & None    & 10 20 30 40 50 60  & 10 20 30 40 50 60       \\
    1   & Up      & 5 5 5 5 5 5        & 15 25 35 45 55 65       \\
    2   & Sub     & 20 30 40 30 30 30  & 20 30 40 50 60 70       \\
    3   & Average & 15 20 25 18 18 18  & 25 35 45 55 65 75       \\
    4   & Paeth   & 5 5 5 5 5 5        & 30 40 50 60 70 80       \\
    \bottomrule
  \end{tabular}
  \caption{Five filters reconstruct a five-row image, one row at a time.}
\end{table}

For the second red value in row two, residual 30 adds the newly reconstructed
left red value 20, yielding 50. In row three the second red predictor is
$\lfloor(25+50)/2\rfloor=37$; residual 18 restores 55. Using residuals as
left neighbors would produce the wrong recurrence.

\texttt{png\_unfilter} and \texttt{png\_unfilter\_msa} receive buffer, exact
count, width and height. Before changing any sample, validate dimensions,
$h(1+3w)$ and every selector. A bad selector in the final row must leave all
rows unchanged. On success, reconstruct in place while retaining selectors;
the same buffer still occupies 35 bytes. The two native views below show the
first sixteen bytes before and after that operation. The Up residuals become
samples, while row boundaries keep their positions.
\begin{figure}[!ht]
  \centering
  \begin{tikzpicture}[alt={A native Memory row contains the first sixteen inflated bytes. An outline marks six Up residual bytes, all 05. Dashed outlines identify the preceding None and Up selectors.}]
    \node[anchor=north west,inner sep=0,font=\small\bfseries] (title) at (0,0) {The Up row still contains residuals};
    \node[anchor=north west,inner sep=0] (a) at ([yshift=-2mm]title.south west) {\MARSRevisionCrop{revision/captures/png/png-inflated.png}{0.99\linewidth}{0.03125}{0.33888889}{0.58203125}{0.60888889}};
    \begin{scope}[shift={(a.south west)},x={($(a.south east)-(a.south west)$)},y={($(a.north west)-(a.south west)$)}]
      \draw[black,line width=.7pt] (.5825,.01) rectangle (.869375,.467);
      \draw[black,dashed,line width=.7pt] (.20,.01) rectangle (.2478125,.467);
      \draw[black,dashed,line width=.7pt] (.5346875,.01) rectangle (.5825,.467);
    \end{scope}
  \end{tikzpicture}
  \caption{Before reconstruction, the six outlined bytes are Up residuals of five. The dashed selectors are zero for the first row and two for the second. The final two pictured bytes begin the next row; they are outside the outlined Up samples.}
  \label{fig:revision-png-native-inflated}
\end{figure}

\begin{figure}[!ht]
  \centering
  \begin{tikzpicture}[alt={The same native Memory interval after reconstruction. Six outlined bytes are 0f 19 23 2d 37 41, representing 15, 25, 35, 45, 55 and 65. Selectors zero and two retain their positions.}]
    \node[anchor=north west,inner sep=0,font=\small\bfseries] (title) at (0,0) {The same bytes now represent RGB samples};
    \node[anchor=north west,inner sep=0] (a) at ([yshift=-2mm]title.south west) {\MARSRevisionCrop{revision/captures/png/png-reconstructed.png}{0.99\linewidth}{0.03125}{0.33888889}{0.58203125}{0.60888889}};
    \begin{scope}[shift={(a.south west)},x={($(a.south east)-(a.south west)$)},y={($(a.north west)-(a.south west)$)}]
      \draw[black,line width=.7pt] (.5825,.01) rectangle (.869375,.467);
      \draw[black,dashed,line width=.7pt] (.20,.01) rectangle (.2478125,.467);
      \draw[black,dashed,line width=.7pt] (.5346875,.01) rectangle (.5825,.467);
    \end{scope}
  \end{tikzpicture}
  \caption{After reconstruction, the outlined samples are decimal 15, 25, 35, 45, 55 and 65. The selectors remain zero and two. The complete buffer retains its thirty-five-byte extent, although its sample bytes have changed meaning.}
  \label{fig:revision-png-native-reconstructed}
\end{figure}


Up has no dependency on a newly reconstructed byte in its current row: add the
previous row's samples with vector byte addition for complete groups, then a
bounded scalar tail. Sub, Average and Paeth depend on reconstructed left
samples. Loading a vector of left residuals does not satisfy that dependency.
The worked MSA route accelerates safe Up groups and packing; selecting MSA
mode does not imply that every stage becomes vector code.

\section{Publish numerical words within active rows}

RGB bytes \texttt{0a 14 1e} represent the first pixel. Its display word is
\texttt{0x000a141e}. A little-endian word store writes \texttt{1e 14 0a 00};
a big-endian store writes \texttt{00 0a 14 1e}. Both store the same numerical
color. Never infer a word's meaning solely from the order of a screenshot's
bytes without knowing the target byte order.

Packing accepts reconstructed rows, exact count, framebuffer and a descriptor
with native width, height, stride in words and capacity in bytes. Require a
word-aligned destination, $w\leq s\leq512$, and capacity at least
\[
  4((h-1)s+w).
\]
For width two, height five and stride seventeen, the selected extent is 280
bytes, but only forty bytes belong to pixels. The intervening sixty-byte gaps
between rows remain caller-owned padding. No trailing gap is needed after the
last active row. Validate this layout before any framebuffer store.

For sixteen pixels, load exactly 48 RGB bytes into three byte vectors and
produce four word vectors. Pixel five straddles the first two input vectors:
its RGB indices are 15, 16 and 17. The second output vector uses selectors
\texttt{14,13,12,128}, \texttt{17,16,15,128}, and so on, creating B,G,R,zero
byte lanes for pixels four through seven. A selector with bit seven set gives
zero. With that retained mask in \texttt{\$w5}, the two input vectors in
\texttt{\$w0/\$w1}, and output base in \texttt{\$t0}:
\Needspace{8\baselineskip}
\MARSCodeLabel{Pack words across the first input-vector boundary}
\begin{minted}{mars-idea}
    move.v $w3, $w5
    vshf.b $w3, $w1, $w0
    st.w $w3, 16($t0)
\end{minted}
The destination initially holds selectors and is overwritten by the shuffle,
so copy a retained mask for each use. Store words with \texttt{st.w} to preserve
native numerical colors in either byte order. Prove a full sixteen-pixel group
exists before the three loads; finish a shorter row with the scalar tail.
Do not borrow bytes from a row selector or padding to make a vector load fit.
\begin{figure}[!ht]
\centering
\begin{tikzpicture}[alt={A native Memory row at the framebuffer base. The first eight bytes encode numeric pixel words 000a141e and 0028323c in little-endian order. The following eight padding bytes remain zero.}]
\node[anchor=north west,inner sep=0,font=\small\bfseries] (title) at (0,0) {Two pixel words precede row padding};
\node[anchor=north west,inner sep=0] (a) at ([yshift=-2mm]title.south west) {\MARSRevisionCrop{revision/captures/png/png-packed.png}{0.99\linewidth}{0.03125}{0.33888889}{0.58203125}{0.60888889}};
\begin{scope}[shift={(a.south west)},x={($(a.south east)-(a.south west)$)},y={($(a.north west)-(a.south west)$)}]
\draw[black,line width=.7pt] (.20,.01) rectangle (.5825,.467);
\draw[black,dashed,line width=.7pt] (.5825,.01) rectangle (.965,.467);
\end{scope}
\end{tikzpicture}
\caption{The solid outline contains the first two numerical pixel words, 000a141e and 0028323c. The dashed bytes are the beginning of caller padding, still zero. With the ordinary stride, the next image row starts 2048 bytes after this one.}
\label{fig:revision-png-native-packed}
\end{figure}


The complete decoder publishes its context results only after collection,
exact inflation, reconstruction and packing succeed. A format or capacity
failure leaves the whole context and framebuffer unchanged, although both
workspaces may have changed. An image with valid chunk CRCs and Adler-32 but
a bad final filter must therefore display no partial result. On success the
complete decoder returns zero with $4wh$ in \texttt{v1}, and publishes width,
height and that active byte count in the context.
\begin{figure}[!ht]
  \centering
  \begin{tikzpicture}[alt={Two native Memory rows show the end of the caller context. The final eight bytes of the first row hold width two; the second row holds height five and active pixel byte count forty. Its first eight pictured bytes are the input row capacity, not an image result.}]
    \node[anchor=north west,inner sep=0,font=\small\bfseries] (title) at (0,0) {Dimensions and count appear after pixel publication};
    \node[anchor=north west,inner sep=0] (a) at ([yshift=-2mm]title.south west) {\MARSRevisionCrop{revision/captures/png/png-context.png}{0.99\linewidth}{0.03125}{0.31555556}{0.58203125}{0.60888889}};
    \begin{scope}[shift={(a.south west)},x={($(a.south east)-(a.south west)$)},y={($(a.north west)-(a.south west)$)}]
      \draw[black,line width=.7pt] (.5825,.3088) rectangle (.965,.6176);
      \draw[black,line width=.7pt] (.20,.005) rectangle (.5825,.303);
      \draw[black,dashed,line width=.7pt] (.5825,.005) rectangle (.965,.303);
    \end{scope}
  \end{tikzpicture}
  \caption{Width two occupies the solid outline at the right of the first row. Height five follows in the next row, then the dashed active-byte result 28 hexadecimal, or forty decimal. The first eight pictured bytes hold the pre-existing row capacity and are outside these three results.}
  \label{fig:revision-png-native-context}
\end{figure}


\section{Run, compare and change the prediction direction}

From the public book repository, generate the trace before launching the worked
caller. The script refuses an existing output. Here \texttt{CLI} is the full
path of the built MARS CLI JAR.
\Needspace{9\baselineskip}
\MARSCodeLabel{Generate the small image and run its complete decoder}
\begin{Verbatim}
python3 packages/033-png/construction/png-fixtures/fixture.py image.png
java -jar "$CLI" --start-at-main revision/labs/files/png-viewer.asm
\end{Verbatim}
The ordinary process working directory must contain \texttt{image.png}.
In IntelliJ, select the caller for MARS Run or Debug and set that working
directory explicitly. The successful console line is \texttt{PNG displayed.}
Inspect the bitmap tool together with the context and rows; the message alone
is not evidence of exact pixels. Set \texttt{vector\_mode} to zero for scalar
or one for MSA. Compare the same file and layout, first for exact state and then
for measured work. File loading, assembly and the complete decode call are
different measurement scopes.

Next generate a fresh seventeen-by-five image with fixed compression and one
with automatic compression. Those names select encoder strategies; automatic
compression does not promise a dynamic block on every input. Record the actual
block types when studying a stream. Change the IDAT split without changing the
joined payload, then verify that decoded samples remain identical. Widths
fifteen, sixteen and seventeen distinguish a scalar tail, one full packing
group, and a group followed by a tail.

The independent \texttt{png-filter-scanline.asm} reverses the row transformation:
from known RGB and a selected predictor, emit the selector and residual bytes.
Compute each residual as $(sample-prediction)\bmod256$. It uses original
current-row neighbors, so a backward traversal is valid when inputs and output
are disjoint. This differs from reconstruction's left-to-right dependency.
The lab specifies the complete interface, tie cases and inverse checks.

\section*{Sources and laboratory connection}
\path{packages/033-png/} retains the exact accepted profile, independent entry
contracts and staged construction material. The
\href{https://www.w3.org/TR/png-3/}{PNG specification},
\href{https://www.rfc-editor.org/rfc/rfc1950}{RFC 1950} and
\href{https://www.rfc-editor.org/rfc/rfc1951}{RFC 1951} define the container,
zlib wrapper and DEFLATE syntax. Local size, publication and preservation
contracts are stated separately in this exercise. Use the written lab to
reason about the boundaries before implementing the complete decoder from its
starter or extending the worked program.
